\section{Is the Gap Knowledge? An 8B LLM Classifier}\label{sec:llm}

The encoder track and the error analysis (\S\ref{sec:analysis}) show that the model, not label noise, is what fails.
Two explanations remain: the encoder may lack \emph{knowledge} that a 0.4B-parameter model has not acquired, or \emph{label meaning}, which a classifier must infer from about 3,400 examples.
A larger pretrained classifier tests the first with nothing else changed.
Before any run we registered three predictions: the gain passes the screening bar; the single-model \Stwo{} mean lies in 0.43--0.50 and \Sone{} rises by +0.02 to +0.05; and the gain falls mostly on the commitment boundary (Explicit $\leftrightarrow$ Implicit, General $\leftrightarrow$ Implicit/Explicit) and more on agreed than on contested items.

\paragraph{One change: the backbone.}
We took the best encoder configuration (sub-question, full question and answer, 16 epochs) and replaced DeBERTa-v3-large with Qwen3-8B-Base \citep{qwen3}, adapted with LoRA \citep{hu2022lora} of rank 16 ($\alpha=32$, dropout 0.05) on all linear projections of the attention and feed-forward blocks.
A new 9-way linear head, trained in full, reads the hidden state of the last token (Figure~\ref{fig:arch}).
The input text and budgets are those of the encoder, with one fixed closing question appended (``How does the answer respond to the sub-question?''), so the token the head reads sits after the whole question and answer in every item.
Sequences are right-padded, and each run first checks that an item's logits are the same alone and inside a padded batch, aborting otherwise; this guards against a head that reads a padding token.
Training ran in bf16 with AdamW at $10^{-4}$, a cosine schedule with 5\% warm-up, 3 epochs and an effective batch of 16; the learning rate and epoch count change because the backbone does.
Everything else is identical: the same 3,103 training rows, the same 345-row train slice (asserted identical in code), plain cross-entropy and the epoch with the best slice macro-F1.
Seeds 0--9 are paired with the DeBERTa runs of the same seeds; Appendix~\ref{sec:app} lists all hyperparameters and the compute.

\begin{table*}[t]
\centering
\small
\setlength{\tabcolsep}{5pt}
\begin{tabular}{@{}lccccc@{}}
\toprule
Dev macro-F1 & Qwen3-8B + LoRA & DeBERTa-v3-large & $\Delta$ & Seeds up & 95\% CI \\
\midrule
Single model, \Stwo{}            & \pmsd{0.476}{0.043} & \pmsd{0.384}{0.030} & +0.092 & 10/10 & \\
Single model, \Sone{}            & \pmsd{0.709}{0.031} & \pmsd{0.614}{0.027} & +0.095 & 10/10 & \\
Single model + LA, \Stwo{}       & \pmsd{0.538}{0.043} & \pmsd{0.389}{0.019} & +0.149 & 10/10 & \\
System, \Stwo{}                  & 0.543 (0.549) & 0.405 (0.412) & +0.136 & & [+0.054, +0.224] \\
System, \Sone{}                  & 0.746 & 0.648 & +0.098 & & \\
System, \Stwo{}, 2-annotator sets & 0.524 & 0.382 & +0.142 & & \\
\midrule
Follow-up (one change)           & Variant & Qwen, 3 epochs & $\Delta$ & Seeds up & 95\% CI \\
\midrule
12 epochs, single model, \Stwo{}$^{\ddagger}$ & \pmsd{0.543}{0.015} & \pmsd{0.462}{0.080} & +0.082$^{\dagger}$ & & \\
All of train, single model, \Stwo{}          & \pmsd{0.492}{0.045} & \pmsd{0.476}{0.043} & +0.017 & 6/10 & \\
All of train, system, \Stwo{}                & 0.575 (0.562) & 0.543 (0.549) & +0.028 & & [$-$0.026, +0.092] \\
\bottomrule
\end{tabular}
\caption{The backbone swap (top) and two follow-ups (bottom) on dev. Single model: mean\,$\pm$\,s.d.\ over seeds, $\Delta$ the paired mean difference. LA: logit adjustment. System: 10-seed ensemble + LA, $\tau$ by nested CV (in parentheses: mean over 10 CV splits); CIs are bootstraps over dev items. 2-annotator sets: the test regime. $^{\ddagger}$Seeds 0--2. $^{\dagger}$Reported only; the 12-epoch schedule was adopted on the train slice.}
\label{tab:llm-main}
\end{table*}

\paragraph{Results at ten paired seeds.}
The 3-seed screen passed the bar (+0.081 on \Stwo{}, 3 of 3 seeds up), so the run was extended to ten seeds (Table~\ref{tab:llm-main}).
Per model, Qwen improves \Stwo{} by +0.092 ($t=7.12$) and \Sone{} by +0.095 ($t=8.48$) on 10 of 10 seeds, more than the whole encoder track gained (0.315 to 0.384).
As a system it reaches 0.543 / 0.746 against 0.405 / 0.648, and the ranking holds on the 2-annotator reference sets (0.524 against 0.382).
All registered size predictions held (the \Sone{} gain exceeded its range); the prediction about \emph{where} the gain falls was refuted (\S\ref{sec:analysis}).
Among published dev results (Figure~\ref{fig:llm}), the system lies above TeleAI's fine-tuned Qwen2.5-7B (0.495) and ChulaNLP's pipeline (0.52), and below TeleAI's multi-call pipeline (0.617) and the human reference (0.684); these come from other protocols, so the comparison locates our system rather than ranking it.

\FigureLLM

\paragraph{Two single-change follow-ups.}
In the 3-epoch runs the slice F1 rose at every epoch and the last epoch was always selected; under a cosine schedule this proves nothing, because the schedule itself favours the last epoch, so only a longer schedule tests undertraining.
With 12 epochs (seeds 0--2) the slice selected epochs 9, 11 and 8 and the slice F1 rose by +0.084 on 3 of 3 seeds, so the rule registered in advance adopted the longer schedule; training loss fell to about 0.02 without the slice F1 declining, and dev \Stwo{} was \pmsd{0.543}{0.015} (reported only).
Training on all 3,448 rows gave +0.017 per model (6 of 10 seeds up), but the system difference has a 95\% interval of [$-$0.026, +0.092] and \Sone{} fell from 0.746 to 0.711; as with the encoder, the 3-seed screen (+0.039) overstated the gain.

\paragraph{Qualitative analysis: what the backbone changes.}
Item by item, the Qwen system moves 65 dev items into the reference set and 36 out of it (194 in-set against 165).
The clearest fixes are answers whose \emph{form} signals the type: Declining to answer (5 fixed, 1 lost), Claims ignorance (5, 0) and Clarification (1, 0).
Asked whether he could comment on concerns for the region, the speaker replied \emph{``No. Saying I'm not going to comment on the matter means I'm not going to comment on the matter.''}; all three annotators chose Declining to answer, DeBERTa chose Claims ignorance and Qwen was right.
The question asked back, \emph{``On Iran?''}, is a unanimous Clarification that DeBERTa called Dodging.
Qwen's remaining errors keep the commitment pattern but move inside the Ambivalent branch: its most frequent error pairs (reference $\rightarrow$ prediction) are Dodging $\rightarrow$ Implicit (17), General $\rightarrow$ Implicit (16) and Implicit $\rightarrow$ Explicit (14), and it over-predicts Deflection (42 predictions against 32 for DeBERTa).
Some losses are literal readings: \emph{``I thought you were going to ask me about the pig''}, in reply to a question about the risk of a wider war, is a unanimous Dodging that Qwen labels Clarification.
(Examples are the shortest-answer item of each category, chosen by rule; the full listing is in the repository.)
